\documentclass[11pt]{article}

\usepackage[T1]{fontenc}
\usepackage{lmodern}
\usepackage{microtype}
\usepackage{amsmath,amssymb,amsthm,mathtools}
\usepackage{aliascnt}
\usepackage[margin=1in]{geometry}
\usepackage[numbers,sort&compress]{natbib}
\usepackage[colorlinks=true,allcolors=blue]{hyperref}

\allowdisplaybreaks
\emergencystretch=2em

\newtheorem{theorem}{Theorem}[section]
\newaliascnt{lemma}{theorem}
\newtheorem{lemma}[lemma]{Lemma}
\aliascntresetthe{lemma}
\newaliascnt{proposition}{theorem}
\newtheorem{proposition}[proposition]{Proposition}
\aliascntresetthe{proposition}
\newaliascnt{corollary}{theorem}
\newtheorem{corollary}[corollary]{Corollary}
\aliascntresetthe{corollary}
\theoremstyle{definition}
\newaliascnt{definition}{theorem}
\newtheorem{definition}[definition]{Definition}
\aliascntresetthe{definition}
\newaliascnt{example}{theorem}
\newtheorem{example}[example]{Example}
\aliascntresetthe{example}
\newaliascnt{remark}{theorem}
\newtheorem{remark}[remark]{Remark}
\aliascntresetthe{remark}
\usepackage[nameinlink,noabbrev]{cleveref}

\newcommand{\E}{\mathsf E}
\newcommand{\dsmall}{\mathsf d}
\newcommand{\D}{\mathsf D}
\newcommand{\Dpm}{\mathsf D_{\pm}}
\newcommand{\Dih}{\operatorname{Dih}}
\newcommand{\F}{\mathcal F}
\newcommand{\supp}{\operatorname{supp}}
\newcommand{\ord}{\operatorname{ord}}
\newcommand{\id}{1_G}

\title{The Gao Constant of Generalized Dihedral Groups\\
with an Abelian \(p\)-Group Kernel, \(p\) Odd}
\author{}
\date{}

\hypersetup{
  pdftitle={The Gao Constant of Generalized Dihedral Groups with an Abelian p-Group Kernel, p Odd},
  pdfsubject={A candidate exact formula for the Gao constant of generalized dihedral groups},
  pdfkeywords={Gao constant, product-one sequence, generalized dihedral group, Davenport constant, weighted zero-sum sequence}
}

\begin{document}
\maketitle

\begin{abstract}
Let \(A\) be a nontrivial finite abelian \(p\)-group for an odd prime
\(p\), and let \(\Dih(A)=A\rtimes_{-1}C_2\) be the generalized
dihedral group in which the nonidentity element of \(C_2\) acts by
inversion. We determine its Gao constant:
\[
  \E(\Dih(A))=2|A|+\D(A)
  =|\Dih(A)|+\dsmall(\Dih(A)).
\]
Thus, if \(A=\bigoplus_{i=1}^r C_{p^{\lambda_i}}\), then
\[
  \E(\Dih(A))
  =2|A|+1+\sum_{i=1}^r\bigl(p^{\lambda_i}-1\bigr).
\]
The lower bound follows from the known small Davenport constant of
\(\Dih(A)\). For the upper bound, we divide according to the number of
reflections in the sequence. The two extreme ranges are handled by
prescribed-length weighted zero-sum results and a plus-minus Davenport
bound. In the remaining range, a weighted structural theorem either
gives the required subsequence directly or forces many rotations into a
proper subgroup. A simultaneous induction on the subgroup order
completes the proof.
\end{abstract}

\noindent\textbf{2020 Mathematics Subject Classification.} 11B75; 20K01.\par
\noindent\textbf{Keywords.} Gao constant; product-one sequence; generalized dihedral group; Davenport constant; weighted zero-sum sequence.

\section{Introduction}\label{sec:introduction}

Let \(G\) be a finite group. Its \emph{Gao constant} \(\E(G)\) is the
least integer \(L\) such that every sequence of length at least \(L\)
over \(G\) contains a subsequence of exactly \(|G|\) terms that can be
ordered to have product one. The small Davenport constant
\(\dsmall(G)\) is the largest length of a sequence over \(G\) having no
nonempty product-one subsequence. The standard padding construction
gives
\[
  \E(G)\ge |G|+\dsmall(G).
\]
For finite abelian groups, Gao proved equality
\cite{Gao1996}. Zhuang and Gao conjectured that the same equality holds
for every finite group \cite{ZhuangGao2005}. It is known for ordinary
dihedral groups; see \cite{ZhuangGao2005,Bass2007,GaoLu2008}. Exact
calculations for other nearby nonabelian families include
\(D_{2n}\times C_2\) \cite{BrocheroLemosMoriyaRibas2023}.

Let \(p\) be an odd prime, let
\[
  A=\bigoplus_{i=1}^{r} C_{p^{\lambda_i}},
  \qquad 1\le \lambda_1\le\cdots\le\lambda_r,
\]
and write
\[
  \Dih(A)=A\rtimes_{-1}C_2,
\]
where the nonidentity element of \(C_2\) acts on \(A\) by inversion.
Our main result is the following.

\begin{theorem}\label{thm:main}
Let \(Q=|A|\). Then
\[
  \E(\Dih(A))=2Q+\D(A)
  =2Q+1+\sum_{i=1}^{r}\bigl(p^{\lambda_i}-1\bigr).
\]
Equivalently,
\[
  \E(\Dih(A))=|\Dih(A)|+\dsmall(\Dih(A)).
\]
\end{theorem}

The equivalence uses the identity
\(\dsmall(\Dih(A))=\D(A)\) of Godara, Joshi, and Mazumdar
\cite{GodaraJoshiMazumdar2026}. Their work also gives the upper bound
\(\E(\Dih(A))\le 3|A|\). When \(A\) is cyclic, this upper bound meets
the standard lower bound and recovers the known dihedral formula. For a
noncyclic abelian \(p\)-group one has \(\D(A)<|A|\), so that bound does
not determine the exact value. More recently, Zhao and Wang proved a
uniform \(3|G|/2\) bound for noncyclic groups whose order is not
divisible by four, together with an equality characterization
\cite{ZhaoWang2027}. Applied here, their result gives a strict upper
bound below \(3|A|\) when \(A\) is noncyclic, but it still does not give
the value in \Cref{thm:main}.

The theorem extends the cyclic-kernel formula to all invariant-factor
decompositions of finite abelian \(p\)-groups with \(p\) odd.

\begin{example}\label{ex:C3square}
Let \(A=C_3^2\), with basis \(e_1,e_2\). Then \(|A|=9\) and
\(\D(A)=5\), so \Cref{thm:main} gives
\[
  \E(\Dih(C_3^2))=23.
\]
The lower bound can be seen explicitly. Take two copies of the rotation
\((e_1,0)\), two copies of \((e_2,0)\), and one reflection. This
five-term sequence is product-one-free: a product-one subsequence has an
even number of reflections, while the four rotations form a zero-sum-free
sequence in \(C_3^2\). Appending seventeen identity rotations gives a
sequence of length \(22\) with no product-one subsequence of length
\(18\).
\end{example}

We outline the upper bound. Let \(S\) have length
\(2|A|+\D(A)\), and let \(a\) be its number of reflections. If
\(a\le 1\), there are enough rotations to apply an unweighted
prescribed-length zero-sum theorem. If \(a\ge \D(A)+1\), pairing terms
within the two cosets reduces the problem to a plus-minus weighted
zero-sum problem in \(A\). In the intermediate range, the structural
alternative in a theorem of Grynkiewicz, Marchan, and Ordaz
\cite{GrynkiewiczMarchanOrdaz2012} either produces the desired
subsequence or shows that many rotations lie in a proper subgroup
\(K<A\). The latter situation is resolved by simultaneous induction on
\(|K|\). The auxiliary induction statement is needed because one branch
translates all rotation vectors; only an all-rotation conclusion can be
pulled back through that translation without changing the product-one
condition.

The proof uses Olson's formula for the Davenport constant of an abelian
\(p\)-group and a plus-minus argument over a field of odd
characteristic. We do not treat mixed-prime kernels or \(2\)-groups in
this paper.

\section{Preliminaries}\label{sec:preliminaries}

A sequence over a set \(X\) is an element of the free abelian monoid
\(\F(X)\). We write it multiplicatively, even when \(X\) is an
additively written group. Thus
\[
  S=x_1\cdot\ldots\cdot x_\ell,
  \qquad |S|=\ell.
\]
Subsequences are selected by positions, so equal values occurring in
different positions remain available separately. We write \(T\mid S\)
when \(T\) is a subsequence and \(ST^{-1}\) for the complementary
sequence. A sequence over a group is \emph{product one} if its terms can
be ordered so that their product is the identity.

For an abelian group \(B\), the Davenport constant \(\D(B)\) is the
least integer \(L\) such that every sequence of length \(L\) over
\(B\) has a nonempty zero-sum subsequence. Thus
\(\D(B)=\dsmall(B)+1\). If \(W\subseteq\mathbb Z\) is nonempty, the
weighted Davenport constant \(\D_W(B)\) is defined analogously, with a
coefficient from \(W\) assigned to each selected term. We write
\[
  \Dpm(B)=\D_{\{-1,1\}}(B).
\]

Throughout the paper, set
\[
  q_i=p^{\lambda_i},\qquad
  Q=|A|=\prod_{i=1}^r q_i,
  \qquad
  D_A=\D(A)=1+\sum_{i=1}^r(q_i-1).
\]
The formula for \(D_A\) is Olson's theorem
\cite{Olson1969}. In particular,
\begin{equation}\label{eq:basic-A-facts}
  D_A\le Q,
  \qquad \exp(A)\mid Q,
  \qquad D_A\ \text{and}\ Q\ \text{are odd}.
\end{equation}
The first inequality follows by expanding
\(Q=\prod_i(1+(q_i-1))\).

We represent elements of \(G=\Dih(A)\) as pairs
\((x,\varepsilon)\in A\times C_2\), with multiplication
\begin{equation}\label{eq:group-law}
  (x,\varepsilon)(y,\delta)
  =\bigl(x+(-1)^\varepsilon y,\,\varepsilon+\delta\bigr).
\end{equation}
Elements \((x,0)\) are called \emph{rotations}, and elements
\((x,1)\) are called \emph{reflections}.

\subsection{Realizing prescribed signs}

The next elementary lemma is the bridge between weighted sums in \(A\)
and product-one subsequences in \(G\).

\begin{lemma}[Realization of prescribed signs]\label{lem:sign-realization}
Suppose a selected sequence contains \(2k>0\) reflections and any
number of rotations. Assign signs from \(\{+1,-1\}\) to the terms so
that exactly \(k\) reflections receive each sign; the signs on the
rotations are arbitrary. Then the terms can be ordered so that their
contributions to the \(A\)-coordinate in \eqref{eq:group-law} have the
assigned signs. Consequently, if the assigned signed vector sum is
zero, the selected sequence is product one.

If a selected sequence contains no reflections, it is product one if
and only if the sum of its rotation vectors is zero.
\end{lemma}

\begin{proof}
Place the positively and negatively signed reflections alternately,
starting with a positive one. In any ordered product, a term contributes
its vector with sign \((-1)^t\), where \(t\) is the number of preceding
reflections. Hence the chosen reflection ordering realizes the prescribed
alternating signs.

The \(2k\) reflections determine \(2k+1\) gaps: one before the first
reflection, one after the last, and one between each consecutive pair.
A rotation inserted in a gap preceded by an even number of reflections
has sign \(+1\); in a gap preceded by an odd number it has sign \(-1\).
Since \(k\ge1\), gaps of both parities occur, and any number of rotations
can be inserted in either type of gap. This realizes all prescribed
rotation signs. The total number of reflections is even, so the
\(C_2\)-coordinate is zero; the \(A\)-coordinate is the prescribed
signed sum.
\end{proof}

\subsection{Prescribed-length weighted sums}

For a sequence \(T\) over a finite abelian group \(B\), a nonempty
weight set \(W\subseteq\mathbb Z\), and an integer \(n\le |T|\), let
\[
  \Sigma_n^W(T)
  =\left\{
      \sum_{j=1}^n w_jt_j:
      t_1,\ldots,t_n\text{ occur in distinct positions of }T,
      \ w_j\in W
    \right\}.
\]
For an integer \(n\), write \(nB=\{nx:x\in B\}\).
We use the following two corollaries of Grynkiewicz, Marchan, and Ordaz
\cite[Corollaries 1.2 and 1.3]{GrynkiewiczMarchanOrdaz2012}.

\begin{theorem}[Grynkiewicz--Marchan--Ordaz]\label{thm:GMO}
Let \(B\) be a finite abelian group, let \(W\subseteq\mathbb Z\) be
nonempty, let \(T\) be a sequence over \(B\), and let \(n\ge |B|\).
If
\[
  |T|\ge n+\D_W(B)-1,
\]
then
\begin{equation}\label{eq:GMO-existence}
  \Sigma_n^W(T)\cap nB\ne\varnothing.
\end{equation}
If, in addition, \(\gcd(W)=1\), then either
\begin{equation}\label{eq:GMO-full}
  \Sigma_n^W(T)=B,
\end{equation}
or there exist a proper subgroup \(H<B\), elements
\(\alpha,\beta\in B\), and a subsequence \(T'\mid T\) such that
\begin{equation}\label{eq:GMO-concentration}
  |T'|\ge |T|-|B/H|+2,
  \qquad
  \supp(T')\subseteq \alpha+H,
\end{equation}
and
\begin{equation}\label{eq:GMO-weight-coset}
  W x\subseteq\beta+H
  \quad\text{for every }x\in\supp(T').
\end{equation}
\end{theorem}

Two specializations will be used repeatedly. For \(W=\{1\}\), one has
\(\D_W(B)=\D(B)\). For \(W=\{-1,1\}\), condition
\eqref{eq:GMO-weight-coset} says that \(x\) and \(-x\) lie in the same
\(H\)-coset for every term of \(T'\). Thus \(2x\in H\); if \(B/H\)
has odd order, this implies \(x\in H\).

The other external input is the following consequence of
\cite[Theorem 1.1]{GodaraJoshiMazumdar2026}:
\begin{equation}\label{eq:small-Davenport-dihedral}
  \dsmall(\Dih(B))=\D(B)
  \qquad\text{for every finite abelian group }B.
\end{equation}

\subsection{The lower bound and identity padding}

\begin{proposition}\label{prop:lower-bound}
For \(G=\Dih(A)\),
\[
  \E(G)\ge 2Q+D_A.
\]
\end{proposition}

\begin{proof}
By \eqref{eq:small-Davenport-dihedral}, there is a product-one-free
sequence \(U\) over \(G\) of length \(D_A\). Append \(2Q-1\) copies
of the identity. Every subsequence of length \(2Q\) uses a nonempty
part of \(U\). If such a subsequence were product one, deleting the
identity terms from a product-one ordering would give a nonempty
product-one subsequence of \(U\), a contradiction.
\end{proof}

\begin{lemma}[Identity padding]\label{lem:identity-padding}
Let \(Y\) be a sequence over \(G\) of length \(2Q+D_A\). If at least
\(D_A\) terms of \(Y\) are the identity rotation, then \(Y\) contains
a product-one subsequence of length \(2Q\).
\end{lemma}

\begin{proof}
Let \(m_0\) be the number of identity rotations. From the sequence of
all nonidentity terms, successively remove nonempty product-one
subsequences until the remainder \(R\) is product-one-free. Let \(P\)
be the union of the removed subsequences. Concatenating one product-one
ordering for each removed block makes \(P\) product one. By
\eqref{eq:small-Davenport-dihedral}, \(|R|\le D_A\), and hence
\[
  |P|=(2Q+D_A-m_0)-|R|\ge 2Q-m_0.
\]
Since \(m_0\ge D_A\), also \(|P|\le 2Q\). Therefore at most \(m_0\)
unused identity terms are needed to extend \(P\) to length \(2Q\).
\end{proof}

\section{A plus-minus Davenport bound}\label{sec:plus-minus}

The plus-minus Davenport constant has been studied systematically; see,
for example, \cite{MarchanOrdazSchmid2014}. We need an estimate expressed
in terms of the ordinary Davenport constant. The estimate below is
known, but we include a short group-algebra proof because this precise
form is used in each signed prescribed-length argument.

\begin{proposition}\label{prop:plus-minus-bound}
Let \(B\) be a nontrivial finite abelian \(p\)-group, where \(p\) is
odd. Then
\[
  \Dpm(B)\le \frac{\D(B)+1}{2}.
\]
\end{proposition}

\begin{proof}
Write
\[
  B=\bigoplus_{i=1}^s C_{p^{\mu_i}},
  \qquad r_i=p^{\mu_i},
  \qquad D_B=\D(B)=1+\sum_{i=1}^s(r_i-1).
\]
Work in the group algebra \(R=\mathbb F_p[B]\), write \([x]\) for the
basis element indexed by \(x\in B\), and let \(I\) be the augmentation
ideal. If \(f_i\) is an invariant-factor generator and
\(y_i=[f_i]-[0]\), then
\[
  y_i^{r_i}=([f_i]-[0])^{r_i}=[r_if_i]-[0]=0.
\]
The resulting homomorphism
\[
  \mathbb F_p[y_1,\ldots,y_s]/(y_1^{r_1},\ldots,y_s^{r_s})
  \longrightarrow R
\]
is surjective, and both sides have dimension \(|B|\). It is therefore
an isomorphism. Every monomial of total degree at least \(D_B\)
vanishes, so
\begin{equation}\label{eq:augmentation-nilpotence}
  I^{D_B}=0.
\end{equation}

The number \(D_B\) is odd. Put \(m=(D_B+1)/2\), and take arbitrary
terms \(x_1,\ldots,x_m\in B\). For every \(j\),
\[
  [x_j]+[-x_j]-2[0]
  =[-x_j]([x_j]-[0])^2\in I^2.
\]
Hence the product of these \(m\) elements lies in
\(I^{2m}=I^{D_B+1}=0\). The coefficient of \([0]\) in its expansion is
\[
  \sum_{\substack{\varepsilon\in\{-1,0,1\}^m\\
                   \sum_j\varepsilon_jx_j=0}}
       (-2)^{|\{j:\varepsilon_j=0\}|}.
\]
If the sequence \(x_1\cdot\ldots\cdot x_m\) had no nonempty
plus-minus zero-sum subsequence, the only contributing vector would be
\(\varepsilon=0\). The displayed coefficient would then be
\((-2)^m\ne0\) in \(\mathbb F_p\), contradicting the vanishing of the
product. Thus every sequence of length \(m\) has a nonempty plus-minus
zero-sum subsequence.
\end{proof}

\section{Reduction by the number of reflections}\label{sec:reflection-count}

Let \(S\) be an arbitrary sequence over \(G=\Dih(A)\) of length
\begin{equation}\label{eq:ambient-length}
  |S|=2Q+D_A.
\end{equation}
Let \(a\) be the number of reflections in \(S\), and put
\[
  b=2Q+D_A-a
\]
for the number of rotations. For a subgroup \(K\le A\) and a sequence
\(Y\) over \(G\), let \(m_K(Y)\) denote the number of rotations
\((x,0)\) of \(Y\) with \(x\in K\).

\subsection{At most one reflection}

\begin{proposition}\label{prop:few-reflections}
If \(a\le1\), then \(S\) contains a product-one subsequence of length
\(2Q\) consisting entirely of rotations.
\end{proposition}

\begin{proof}
There are at least
\[
  b\ge 2Q+D_A-1=2Q+\D(A)-1
\]
rotations. Apply \Cref{thm:GMO}, in its existence form
\eqref{eq:GMO-existence}, to the rotation sequence over \(A\), with
weight set \(\{1\}\) and target \(n=2Q\). The target satisfies
\(2Q\ge |A|\), and the displayed inequality is the required length
hypothesis. Since \(\exp(A)\mid Q\), one has \(2QA=\{0\}\). Thus
there are exactly \(2Q\) rotations with ordinary sum zero.
\end{proof}

\subsection{Many reflections}

\begin{proposition}\label{prop:many-reflections}
If \(a\ge D_A+1\), then \(S\) contains a product-one subsequence of
length \(2Q\).
\end{proposition}

\begin{proof}
Pair reflections with reflections and rotations with rotations. Since
\(2Q+D_A\) is odd, exactly one term is left unpaired, and the number of
pairs is
\[
  \left\lfloor\frac a2\right\rfloor
  +\left\lfloor\frac b2\right\rfloor
  =Q+\frac{D_A-1}{2}.
\]
Assign to a reflection pair \((x,1)\cdot(y,1)\) the label \(x-y\), and
to a rotation pair \((u,0)\cdot(v,0)\) the label \(u+v\). By
\Cref{prop:plus-minus-bound}, the sequence of pair labels has length at
least
\[
  Q+\Dpm(A)-1.
\]
Apply \eqref{eq:GMO-existence} in \(A\), with weight set
\(\{-1,1\}\) and target \(Q=|A|\). Since \(QA=\{0\}\), we obtain
exactly \(Q\) pairs and a sign for each selected pair such that the
signed sum of their labels is zero.

At most \(\lfloor b/2\rfloor\le Q-1\) of all pairs are rotation pairs,
because \(a\ge D_A+1\). Hence at least one selected pair is a
reflection pair. For a reflection pair with selected weight
\(\varepsilon\), assign signs \(\varepsilon\) and \(-\varepsilon\)
to its two reflections. For a rotation pair, assign the selected weight
to both rotations. The selected reflection signs are balanced, and the
total signed vector sum is zero. By \Cref{lem:sign-realization}, the
\(2Q\) selected terms have a product-one ordering.
\end{proof}

\subsection{The intermediate range}

Assume now that
\begin{equation}\label{eq:middle-range}
  2\le a\le D_A.
\end{equation}
Let \(e\) be the largest even integer not exceeding \(a\), and put
\begin{equation}\label{eq:middle-target}
  \ell=2Q-e.
\end{equation}
Then \(e>0\), \(\ell\ge Q\), and \(\ell\le b\). Moreover,
\begin{equation}\label{eq:middle-surplus}
  b-\ell=
  \begin{cases}
    D_A,& a\text{ even},\\
    D_A-1,& a\text{ odd},
  \end{cases}
  \qquad
  b-\ell\ge \Dpm(A)-1.
\end{equation}
Indeed, \(e\le D_A\le Q\), and the final inequality follows from
\Cref{prop:plus-minus-bound}.

\begin{proposition}[Intermediate-range reduction]
\label{prop:middle-reduction}
Under \eqref{eq:middle-range}, either \(S\) contains a product-one
subsequence of length \(2Q\), or there is a nonzero proper subgroup
\(K<A\) such that
\begin{equation}\label{eq:concentration-K}
  m_K(S)\ge b-[A:K]+2.
\end{equation}
\end{proposition}

\begin{proof}
Apply the structural part of \Cref{thm:GMO} to the sequence of all
\(b\) rotation vectors, with ambient group \(A\), weight set
\(\{-1,1\}\), and target \(\ell\). The target condition
\(\ell\ge|A|\) and the length surplus are exactly
\eqref{eq:middle-surplus}.

If \(\Sigma_\ell^{\{-1,1\}}(S_{\rm rot})=A\), select any \(e\)
reflections and assign half of them each sign. Let their signed vector
sum be \(t\in A\). Fullness supplies \(\ell\) rotations with a signed
sum of \(-t\). The two selections are disjoint, and their total length
is \(e+\ell=2Q\). The sign assignment on the reflections is balanced,
so \Cref{lem:sign-realization} gives a product-one ordering.

Otherwise, \Cref{thm:GMO} gives a proper subgroup \(K<A\) and a
rotation subsequence of length at least \(b-[A:K]+2\) satisfying
\eqref{eq:GMO-weight-coset}. For every one of its values \(x\), the
elements \(x\) and \(-x\) lie in a common \(K\)-coset. Hence
\(2x\in K\). Since \(A/K\) has odd order, multiplication by two is
injective on \(A/K\), and therefore \(x\in K\). This proves
\eqref{eq:concentration-K}.

It remains to exclude \(K=0\) from the second alternative. In that
case at least
\[
  b-Q+2=Q+D_A-a+2\ge D_A
\]
rotations are the identity, so \Cref{lem:identity-padding} gives the
required subsequence directly.
\end{proof}

\section{Subgroup descent in the intermediate range}
\label{sec:subgroup-descent}

Fix, for this section, a sequence \(S\) satisfying
\eqref{eq:ambient-length} and \eqref{eq:middle-range}; retain its
reflection and rotation counts \(a,b\). For a sequence \(Y\) over
\(G\) and a subgroup \(K\le A\), let \(T_K(Y)\) be the sequence over
\(\Dih(A/K)\) obtained by taking all reflections of \(Y\), together
with all rotations whose vectors lie outside \(K\), and projecting their
vectors modulo \(K\). Rotations with vectors in \(K\) are omitted.
Selections in \(T_K(Y)\) always refer to the corresponding positions of
\(Y\).

We prove two statements simultaneously. For every proper subgroup
\(K<A\), let \(\mathcal P_S(K)\) denote the assertion
\[
  m_K(S)\ge b-[A:K]+2
  \quad\Longrightarrow\quad
  S\text{ has a product-one subsequence of length }2Q.
\]
Let \(\mathcal Q(K)\) denote the following universal assertion: every
sequence \(Y\) over \(G\) of length \(2Q+D_A\), with the same counts
\(a,b\), has an all-rotation product-one subsequence of length \(2Q\)
provided that
\begin{enumerate}
\item \(m_K(Y)\ge b-[A:K]+2\), and
\item \(T_K(Y)\) has no product-one subsequence containing a reflection.
\end{enumerate}
The all-rotation conclusion is essential because the proof of
\(\mathcal Q(K)\) translates rotation vectors before descending to a
smaller subgroup.

\begin{theorem}[Subgroup descent]\label{thm:subgroup-descent}
For every proper subgroup \(K<A\), both \(\mathcal P_S(K)\) and
\(\mathcal Q(K)\) hold.
\end{theorem}

\subsection{Three subgroup estimates}

\begin{lemma}[Davenport inequality]\label{lem:Davenport-inequality}
For every subgroup \(K\le A\),
\begin{equation}\label{eq:Davenport-inequality}
  D_A\ge \D(K)+\D(A/K)-1.
\end{equation}
\end{lemma}

\begin{proof}
Take a zero-sum-free sequence in \(K\) of length \(\D(K)-1\) and a
zero-sum-free sequence in \(A/K\) of length \(\D(A/K)-1\), and lift
the latter arbitrarily to \(A\). A zero-sum subsequence of the
concatenation would project to a zero-sum subsequence of the quotient
sequence. It therefore uses no lifted quotient term, after which it
would be a zero-sum subsequence of the sequence in \(K\). Thus the
concatenation is zero-sum-free, giving \eqref{eq:Davenport-inequality}.
\end{proof}

\begin{lemma}[Target-length estimate]\label{lem:target-length}
Let \(0<K<A\), put \(I=[A:K]\) and \(h=|K|\), so that \(Q=Ih\).
Then
\begin{equation}\label{eq:target-length}
  2Q-D_A-I+2\ge Q-I+2\ge h.
\end{equation}
\end{lemma}

\begin{proof}
The first inequality is \(D_A\le Q\). For the second,
\[
  Q-I+2-h=Ih-I-h+2=(I-1)(h-1)+1>0.
\]
\end{proof}

\begin{lemma}[Transitivity of concentration]\label{lem:concentration-transitivity}
Let \(H<K<A\), put \(I=[A:K]\) and \(J=[K:H]\), and let \(Y\) be a
sequence with
\[
  m_K(Y)\ge b-I+2.
\]
If at least \(m_K(Y)-J+2\) of the rotations counted by \(m_K(Y)\)
actually lie in \(H\), then
\[
  m_H(Y)\ge b-[A:H]+2.
\]
\end{lemma}

\begin{proof}
Since \([A:H]=IJ\),
\[
  m_H(Y)\ge b-I-J+4\ge b-IJ+2,
\]
because \(IJ-I-J+2=(I-1)(J-1)+1>0\).
\end{proof}

\begin{lemma}[Correction of a quotient defect]\label{lem:defect-correction}
Let \(K\le A\). Suppose a selected sequence \(U\) has a positive even
number of reflections and admits an ordering whose product has
\(A\)-coordinate \(z\in K\). Let \(C_0\) be a disjoint rotation
sequence, with prescribed signs, whose signed vector sum is \(-z\).
Then \(UC_0\) is product one.
\end{lemma}

\begin{proof}
Read from the fixed ordering of \(U\) the sign of each vector. The
reflection signs are balanced, because they alternate and their number
is positive and even. Together with the prescribed signs on \(C_0\),
the total signed vector sum is zero. Apply
\Cref{lem:sign-realization} to the entire selected sequence.
\end{proof}

\subsection{A reflection-containing quotient block}
\label{subsec:reflection-quotient}

Let \(0<K<A\), suppose the hypothesis of \(\mathcal P_S(K)\) holds,
and let \(C\) be the sequence of the \(M=m_K(S)\) rotations of \(S\)
whose vectors lie in \(K\). Let \(B\) be the sequence of the remaining
\(c=b-M\) rotations, and put \(T=T_K(S)\).

\begin{proposition}\label{prop:reflection-quotient}
If \(T\) has a product-one subsequence containing a reflection, then
either \(S\) has a product-one subsequence of length \(2Q\), or there
is a strict subgroup \(H<K\) such that
\[
  m_H(S)\ge b-[A:H]+2.
\]
\end{proposition}

\begin{proof}
Start with a reflection-containing product-one subsequence of \(T\).
Continue removing nonempty product-one subsequences from the current
remainder until the final remainder \(R\) is product-one-free. Let
\(U\) be the union of all removed blocks. Concatenating one product-one
ordering for each block makes \(U\) product one in \(\Dih(A/K)\), and
\[
  |U|=a+c-|R|.
\]
The sequence \(U\) contains a positive even number of reflections. By
\eqref{eq:small-Davenport-dihedral}, applied to \(A/K\),
\[
  |R|\le \D(A/K).
\]
Set
\[
  \tau=D_A-|R|.
\]
By \Cref{lem:Davenport-inequality} and
\Cref{prop:plus-minus-bound},
\begin{equation}\label{eq:reflection-surplus}
  \tau\ge \D(K)-1\ge \Dpm(K)-1.
\end{equation}

Lift the fixed quotient ordering of \(U\) to the corresponding terms of
\(S\). Its product has \(C_2\)-coordinate zero and is trivial modulo
\(K\), so its \(A\)-coordinate is a defect \(z\in K\). Put
\[
  m=M-\tau=M-D_A+|R|.
\]
The concentration hypothesis and \Cref{lem:target-length} give
\begin{align*}
  m&\ge M-D_A\\
   &\ge b-[A:K]+2-D_A\\
   &=2Q-a-[A:K]+2\\
   &\ge 2Q-D_A-[A:K]+2\\
   &\ge |K|.
\end{align*}
Equation \eqref{eq:reflection-surplus} gives
\(M\ge m+\Dpm(K)-1\). Thus the structural part of
\Cref{thm:GMO} applies to \(C\) in the ambient group \(K\), with
weight set \(\{-1,1\}\) and target \(m\).

If \(\Sigma_m^{\{-1,1\}}(C)=K\), choose \(m\) rotations of \(C\),
with signs, whose signed sum is \(-z\). Their positions are disjoint
from those of \(U\), because \(T\) contains no rotation from \(K\).
By \Cref{lem:defect-correction}, their union with the lift of \(U\) is
product one. Its length is
\begin{align*}
  |U|+m
  &=a+c-|R|+M-D_A+|R|\\
  &=a+b-D_A=2Q.
\end{align*}

If the spectrum is not all of \(K\), the structural alternative gives
a strict subgroup \(H<K\) and at least \(M-[K:H]+2\) terms of \(C\)
satisfying \eqref{eq:GMO-weight-coset}. For each such vector \(x\), the
elements \(x\) and \(-x\) lie in a common \(H\)-coset, so \(2x\in H\).
The quotient \(K/H\) has odd order, and hence \(x\in H\). Now
\Cref{lem:concentration-transitivity} gives
\(m_H(S)\ge b-[A:H]+2\).
\end{proof}

\subsection{The rotation-only quotient case}
\label{subsec:rotation-quotient}

For \(\alpha\in A\) and a sequence \(Y\) over \(G\), define
\(\tau_\alpha(Y)\) by replacing every rotation \((v,0)\) with
\((v-\alpha,0)\), leaving every reflection unchanged. This is a
transformation of the values attached to the positions; it is not a
group homomorphism.

\begin{lemma}[Persistence of the quotient obstruction]
\label{lem:quotient-obstruction}
Let \(H<K\le A\) and \(\alpha\in K\). If \(T_K(Y)\) has no
product-one subsequence containing a reflection, then
\(T_H(\tau_\alpha(Y))\) has no such subsequence.
\end{lemma}

\begin{proof}
Suppose that a reflection-containing product-one subsequence existed in
\(T_H(\tau_\alpha(Y))\), and fix a product-one ordering. Project this
ordered block further modulo \(K\). A selected rotation whose original
vector \(v\) lies in \(K\) has translated vector \(v-\alpha\in K\), so
it becomes the identity modulo \(K\); delete all such identity rotations
from the ordering. A selected rotation whose original vector lies
outside \(K\) has the same image modulo \(K\) before and after the
translation. Reflections are unchanged. The remaining ordered block is
a reflection-containing product-one subsequence of \(T_K(Y)\), a
contradiction.
\end{proof}

\begin{proposition}\label{prop:rotation-quotient}
Let \(0<K<A\). Suppose that \(Y\) has length \(2Q+D_A\), has the same
counts \(a,b\) as \(S\), satisfies
\[
  m_K(Y)\ge b-[A:K]+2,
\]
and \(T_K(Y)\) has no product-one subsequence containing a reflection.
Then either
\begin{enumerate}
\item \(Y\) has an all-rotation product-one subsequence of length
      \(2Q\), or
\item there are \(H<K\) and \(\alpha\in K\) such that
      \(\tau_\alpha(Y)\) satisfies the hypotheses of \(\mathcal Q(H)\).
\end{enumerate}
\end{proposition}

\begin{proof}
Let \(C\) be the sequence of the \(M=m_K(Y)\) rotations in \(K\), and
let \(B\) be the sequence of the remaining \(c=b-M\) rotations. In the
quotient \(A/K\), successively remove nonempty zero-sum subsequences
from the image of \(B\) until a zero-sum-free remainder \(B_0\)
remains. Let \(B'=BB_0^{-1}\), using the corresponding original
positions, and put \(r_0=|B_0|\). The vector sum of \(B'\) is an
element
\begin{equation}\label{eq:rotation-defect}
  z=\sigma(B')\in K.
\end{equation}

The quotient sequence formed by all \(a\) reflections and the terms of
\(B_0\) is product-one-free. Indeed, a rotation-only product-one
subsequence would give a nonempty zero-sum subsequence of \(B_0\), and
a product-one subsequence containing a reflection is excluded by the
hypothesis on \(T_K(Y)\). Hence, by
\eqref{eq:small-Davenport-dihedral},
\begin{equation}\label{eq:rotation-remainder}
  a+r_0\le \D(A/K).
\end{equation}
Set
\[
  d_0=D_A-a-r_0.
\]
Using \Cref{lem:Davenport-inequality} and
\eqref{eq:rotation-remainder},
\begin{equation}\label{eq:rotation-surplus}
  d_0\ge D_A-\D(A/K)\ge \D(K)-1.
\end{equation}

Put \(m=M-d_0\). The concentration hypothesis gives
\begin{align*}
  m&\ge b-[A:K]+2-(D_A-a-r_0)\\
   &=2Q-[A:K]+r_0+2\\
   &\ge 2Q-[A:K]+2\ge |K|,
\end{align*}
where the last inequality follows from
\Cref{lem:target-length}. Equation \eqref{eq:rotation-surplus} gives
\(M\ge m+\D(K)-1\). Apply the structural part of
\Cref{thm:GMO} to \(C\) in \(K\), with weight set \(\{1\}\) and
target \(m\).

Suppose first that \(\Sigma_m^{\{1\}}(C)=K\). Positional
complementation inside \(C\) gives
\[
  \Sigma_{d_0}^{\{1\}}(C)
  =\sigma(C)-\Sigma_m^{\{1\}}(C)=K.
\]
Choose a \(d_0\)-term subsequence \(D_0\mid C\) with
\(\sigma(D_0)=\sigma(C)+z\). The rotation sequences
\(CD_0^{-1}\) and \(B'\) have disjoint positions and total vector sum
\[
  \sigma(C)-\sigma(D_0)+\sigma(B')=-z+z=0.
\]
Their total length is
\begin{align*}
  |CD_0^{-1}|+|B'|
  &=M-d_0+c-r_0\\
  &=b-(D_A-a-r_0)-r_0=2Q.
\end{align*}
This proves the first alternative.

If the spectrum is not all of \(K\), the structural alternative gives
a strict subgroup \(H<K\), an element \(\alpha\in K\), and at least
\(M-[K:H]+2\) terms of \(C\) whose vectors lie in \(\alpha+H\).
After applying \(\tau_\alpha\), these rotations lie in \(H\), and
\Cref{lem:concentration-transitivity} gives
\[
  m_H(\tau_\alpha(Y))\ge b-[A:H]+2.
\]
The translation preserves the length and the counts \(a,b\), while
\Cref{lem:quotient-obstruction} preserves the absence of a
reflection-containing quotient block. Thus \(\tau_\alpha(Y)\) satisfies
the hypotheses of \(\mathcal Q(H)\).
\end{proof}

\begin{lemma}[Pullback through a rotation translation]
\label{lem:translation-pullback}
If the positions of an all-rotation subsequence of length \(2Q\) form a
product-one sequence in \(\tau_\alpha(Y)\), then the same positions form
a product-one sequence in \(Y\).
\end{lemma}

\begin{proof}
Pulling the selected rotations back adds \(\alpha\) to each vector, so
the total vector sum changes by \(2Q\alpha=0\), using
\(\exp(A)\mid Q\). There are no reflections, so no alternating signs
occur.
\end{proof}

\subsection{Simultaneous induction}

\begin{proof}[Proof of \Cref{thm:subgroup-descent}]
We use strong induction on the order of the subgroup.

For \(K=0\), the hypothesis of \(\mathcal P_S(0)\) says that the
number \(m_0\) of identity rotations satisfies
\[
  m_0\ge b-Q+2=Q+D_A-a+2\ge D_A.
\]
Thus \(\mathcal P_S(0)\) follows from
\Cref{lem:identity-padding}.

To prove \(\mathcal Q(0)\), let \(Y\) satisfy its hypotheses, and let
\(m_0\) again be the number of identity rotations. The sequence
\(T_0(Y)\) consists of all reflections and all nonidentity rotations.
Successively remove nonempty product-one subsequences from \(T_0(Y)\)
until a product-one-free remainder remains, and let \(P\) be the union
of the removed blocks. The quotient obstruction in the definition of
\(\mathcal Q(0)\) forces every removed block to consist only of
rotations. Hence \(P\) is an all-rotation product-one sequence. The
terminal remainder has length at most \(D_A\), so
\[
  |P|\ge (2Q+D_A-m_0)-D_A=2Q-m_0.
\]
Moreover, the concentration hypothesis gives
\(m_0\ge b-Q+2\), and therefore
\[
  |P|\le b-m_0\le Q-2<2Q.
\]
There are enough unused identity rotations to extend \(P\) to an
all-rotation product-one subsequence of length \(2Q\). This proves
\(\mathcal Q(0)\).

Now let \(0<K<A\), and assume both statements are known for every
strict subgroup of \(K\). To prove \(\mathcal P_S(K)\), consider
\(T_K(S)\). If it has a product-one subsequence containing a reflection,
\Cref{prop:reflection-quotient} either gives the required subsequence
immediately or gives a strict subgroup \(H<K\) satisfying the hypothesis
of \(\mathcal P_S(H)\). The induction hypothesis applies in the latter
case. If \(T_K(S)\) has no such subsequence,
\Cref{prop:rotation-quotient} either gives an all-rotation subsequence of
length \(2Q\) directly or gives \(H<K\) and \(\alpha\in K\) such that
\(\tau_\alpha(S)\) satisfies the hypotheses of \(\mathcal Q(H)\). The
induction hypothesis gives the required all-rotation subsequence after
translation, and \Cref{lem:translation-pullback} transfers it back to
\(S\).

Finally, let \(Y\) satisfy the hypotheses of \(\mathcal Q(K)\). Its
quotient obstruction already rules out the reflection-containing case.
Apply \Cref{prop:rotation-quotient}. The first alternative is the desired
conclusion. In the second, the translated sequence satisfies
\(\mathcal Q(H)\) for a strict subgroup \(H<K\); apply the induction
hypothesis and then \Cref{lem:translation-pullback}.
\end{proof}

\section{Completion of the proof}\label{sec:completion}

\begin{proof}[Proof of \Cref{thm:main}]
The lower bound is \Cref{prop:lower-bound}. For the upper bound, let
\(S\) be a sequence of length \(2Q+D_A\), and let \(a\) be its number
of reflections. If \(a\le1\), apply \Cref{prop:few-reflections}. If
\(a\ge D_A+1\), apply \Cref{prop:many-reflections}. In the remaining
range \(2\le a\le D_A\), \Cref{prop:middle-reduction} either gives the
required subsequence directly or yields a nonzero proper subgroup
\(K<A\) satisfying the hypothesis of \(\mathcal P_S(K)\).
\Cref{thm:subgroup-descent} then gives a product-one subsequence of
length \(2Q=|G|\). Hence \(\E(G)\le2Q+D_A\), and equality follows.
\end{proof}

\begin{corollary}\label{cor:homocyclic}
If \(A=C_{p^k}^r\), where \(p\) is odd, then
\[
  \E(\Dih(A))=2p^{kr}+r(p^k-1)+1.
\]
In particular,
\[
  \E(\Dih(C_p^r))=2p^r+r(p-1)+1.
\]
\end{corollary}

\begin{proof}
Substitute \(|A|=p^{kr}\) and
\(\D(C_{p^k}^r)=1+r(p^k-1)\) in \Cref{thm:main}.
\end{proof}


\bibliographystyle{plainnat}
\bibliography{references}

\end{document}
